\section{Dartmouth Workshop}

\par
In the early 1950s, there were multiple names given for the field of "thinking
machines":

\begin{itemize}
	\item Cybernetics
	\item Automata theory
	\item Information processing
	\item Thinking machines
	\item Statistical computing
	\item Computational statistics
	\item Predictive analysis
\end{itemize}

These different names given to the field reflected diverse perspectives on how
the field was being approached.

\par
In 1955,
\href{https://en.wikipedia.org/wiki/John_McCarthy_(computer_scientist)}{John
	McCarthy} proposed a summer research project in 1956 at Dartmouth
using the term "Artificial Intelligence" (AI) to describe the field in a
neutral way. During this workshop, several aspects of AI were discussed:

\begin{itemize}
	\item Automatic computers
	\item How can a computer be programmed to use a language
	\item Neuron nets
	\item Theory of the size of a calculation
	\item Self-improvement
	\item Abstractions
	\item Randomness and creativity
\end{itemize}

\par
Focusing on these problems during the summer to attempt to make a significant
advance in AI\cite{McCarthy_Minsky_Rochester_Shannon_2006}. The attendees of
this conference would later be known as the "founding fathers of AI".
